\documentclass[12pt]{article}
\usepackage[utf8]{inputenc}
\usepackage[T1]{fontenc}
\usepackage{hyperref}
\title{Thawa Language Resources for K-12 Curriculum Development}
\author{Compiled by Claude Sonnet 5.5 (claude-sonnet-5-5) for Andr\'as Kornai}
\date{}
\begin{document}
\maketitle

\section{Introduction}
Thawa (ISO 639-3 xtv; Glottocode sout2771; also spelt Thaua, Dhawa, Thauaira) is an Aboriginal language of the far South Coast of New South Wales, Australia. It is associated with the Thaua people around the Twofold Bay area. Sources disagree on its relationship to neighbouring languages: some treat it as a dialect of Southern Coastal Yuin. AUSTLANG notes that the data for it are limited to wordlists. Locations given by Mathews and Tindale conflict with those given by Eades (1976) and Clark (1996).

\section{Primary Educational Resources}
In 2015, local Yuin people worked with Tathra Public School to create an app as a teaching aid for Thawa and the neighbouring Dhurga language, according to Wikipedia. I did not find the app itself or any other textbook or primer.

\section{Linguistic and Reference Materials}
\begin{itemize}
\item Wikipedia, Thawa language: \url{https://en.wikipedia.org/wiki/Thawa_language}
\item AUSTLANG (AIATSIS) entry S52: \url{https://collection.aiatsis.gov.au/austlang/language/S52}
\item Wikipedia, Tharawal languages (related-language context): \url{https://en.wikipedia.org/wiki/Tharawal_languages}
\item Glottolog reference entry: \url{https://glottolog.org/resource/reference/id/38921}
\end{itemize}
The main documentation is John Liston's 1993 wordlist, \emph{The Thaua language: An Aboriginal pictionary from the N.S.W. far South Coast}, as cited by Wikipedia. I did not find a copy online.

\section{Religious and Community Resources}
No Scripture translation or community organisation specific to Thawa was found in the search.

\section{Digital and Audio Resources}
No audio or video for Thawa was found. The Tathra school app mentioned above is the only digital resource reported.

\section{Curriculum Development Context}
Thawa is treated as nearly extinct or extinct, with documentation limited to wordlists. Any heritage curriculum would depend on Yuin community leadership and on AIATSIS and the wordlist material, possibly alongside Dhurga resources. Such a curriculum would have to rely on reclamation work and should be developed with the community.
\end{document}